\subsection{Esprimere le osservabili come operatori}
Il sistema si presenta come combinazione di stati \(\ket{e_i}\) e quando misuro ottengo un \(\lambda_i\) reale associato a \(\ket{e_i}\). Chiamo \(O\) questa osservabile. Allora
\begin{align*}
\text{valore medio} \quad \braket{O}_\psi &= \sum_i \lambda_i P_i = \sum_i \lambda_i \left| \braket{e_i|\psi} \right|^2 \\
&\stackrel{\text{def}}{=} \braket{\psi|O|\psi}
\end{align*}
\[
\text{operatore} \quad O = \sum_i \lambda_i \ket{e_i} \bra{e_i}
\]
\[
O \ket{\psi} = \sum_i \lambda_i \ket{e_i} \braket{e_i|\psi}
\]
\begin{align*}
\braket{O}_\psi &= \sum_{ij} \braket{\psi|e_i} \braket{e_i|O|e_j} \braket{e_j|\psi} = \\
&= \sum_{ijm} \braket{e_i|\psi} \underbrace{\braket{e_i|\lambda_m|e_m}}_{\lambda_m \delta_{im}} \underbrace{\braket{e_m|e_j}}_{\delta_{mj}} \braket{e_j|\psi} = \\
&= \sum_i \lambda_i \braket{\psi|e_i} \braket{e_i|\psi} = \\
&= \sum_i \lambda_i \left| \braket{\psi|e_i} \right|^2
\end{align*}
\[
\braket{O}_\psi = \sum_i \lambda_i \left| \braket{\psi|e_i} \right|^2 \qquad \braket{O}_\psi = \braket{\psi|O|\psi}
\]
Inoltre, la matrice \(O\) è diagonale
\[
O = \begin{pmatrix}
\lambda_1 & 0 \\
0 & \ddots
\end{pmatrix}
\]
In particolare
\[
\lambda_i \text{ è un autovalore reale}
\]
\[
\ket{e_i} \text{ è un autostato}
\]
Quindi ogni osservabile è un operatore tale che ammette autovalori reali e una BON di autostati e viceversa.

\subsubsection{Esempio: operatore come matrice}
\[
O = \ket{\psi_1}\bra{\psi_2} + \ket{\psi_2}\bra{\psi_1}
\]
\[
\ket{\psi_1} = \frac{1}{\sqrt{2}} \left( \ket{+} + i \ket{-} \right) \qquad \ket{\psi_2} = \frac{1}{\sqrt{2}} \left( \ket{+} - i \ket{-} \right)
\]
Poiché \(\braket{\pm|\pm} = 1\) e \(\braket{\pm|\mp} = 0\), ho la BON e
\begin{align*}
O_{11} &= \braket{+|\frac{1}{\sqrt{2}} \left( \ket{+} + i \ket{-} \right) \frac{1}{\sqrt{2}} \left( \bra{+} + i \bra{-} \right)|+} + \\
&+ \braket{+|\frac{1}{\sqrt{2}} \left( \ket{+} - i \ket{-} \right) \frac{1}{\sqrt{2}} \left( \bra{+} - i \bra{-} \right)|+} = \\
&= \frac{1}{2} + \frac{1}{2} = 1
\end{align*}
\[
O_{12} = -\frac{1}{2} + \frac{1}{2} = 0 \qquad O_{21} = \frac{1}{2} - \frac{1}{2} = 0 \qquad O_{22} = -1
\]
\[
O = \begin{pmatrix}
1 & 0 \\
0 & -1
\end{pmatrix}
\]

\subsubsection{Esempio: proiezione di uno stato}
\[
\ket{\psi} = \frac{1}{2} \ket{1} + \frac{1}{\sqrt{2}} \ket{2} + \frac{1}{\sqrt{2}} \ket{3}
\]
Normalizzo
\[
P_i = \frac{|a_i|^2}{\sum_{k=1}^3 |a_k|^2} = \begin{cases}
1/5 \\
2/5 \\
2/5
\end{cases}
\]
e ho le singole probabilità. La probabilità che non si trovi in \(\ket{2}\) è \(1 - P_2 = 3/5\).

\bigskip
\noindent
Potrei definire un'osservabile
\[
O = \lambda \ket{1}\bra{1} + \lambda \ket{3}\bra{3} + \lambda' \ket{2}\bra{2}
\]
e se misuro che il sistema si trova in \(\ket{2}\) so già, altrimenti se non si trova in \(\ket{2}\) è nello stato
\[
\ket{\psi'} = \frac{1}{2} \ket{1} + \frac{1}{\sqrt{2}} \ket{3}
\]
cioè proietto \(\ket{\psi}\) attraverso un operatore \(P\) tale che \(P^2 = P\) attraverso
\[
P_{\ket{S}} = \ket{S}\bra{S}, \qquad \text{non } S \rightarrow P = \mathbbm{1} - \ket{S}\bra{S}
\]

\bigskip
\noindent
Se si trova in \(\ket{2}\) e voglio calcolare la probabiltià con
\[
\ket{\varphi} = \frac{1}{\sqrt{2}} \left( \ket{2} - \ket{3} \right)
\]
\[
\left| \braket{2|\varphi} \right|^2 = \frac{1}{2} = P
\]
Inoltre,
\[
\left| \braket{\varphi|\psi} \right|^2 = 0
\]
cioè \(\ket{\varphi} \perp \ket{\psi}\) ed è l'informazione sul sistema \(\rightarrow\) il proiettore lascia invariato lo \(\ket{\varphi}\), infatti
\[
P = \mathbbm{1} - \ket{\varphi}\bra{\varphi}, \qquad P \ket{\psi} = \ket{\psi}
\]
e quindi la probabilità che il sistema sia in \(\ket{2}\).

\bigskip
\noindent
Quindi il proiettore
\[
P_\varphi = \ket{\varphi}\bra{\varphi}
\]